\documentclass[a4paper,11pt]{article}
\usepackage[T2A]{fontenc}
\usepackage[utf8]{inputenc}
\usepackage{PTSerif}
\usepackage[russian]{babel}
\usepackage{amsmath,amssymb,mathtools}
\usepackage{geometry}
\geometry{left=23mm,right=23mm,top=22mm,bottom=23mm}
\usepackage{microtype}
\usepackage{booktabs,tabularx,array}
\usepackage{enumitem}
\usepackage{tikz}
\usetikzlibrary{arrows.meta,positioning,shapes.geometric,calc}
\usepackage{xcolor}
\usepackage{fancyhdr}
\usepackage{hyperref}
\hypersetup{hidelinks,pdftitle={Теория вероятностей. Независимые испытания и формула Бернулли},pdfauthor={Лёвин Артём Александрович}}
\definecolor{accent}{HTML}{007C80}
\definecolor{darkaccent}{HTML}{005A5E}
\definecolor{textgray}{HTML}{353535}
\color{textgray}
\setlength{\parindent}{0pt}
\setlength{\parskip}{0.50em}
\setlength{\headheight}{23pt}
\setlength{\emergencystretch}{1.5em}
\renewcommand{\arraystretch}{1.27}
\setlist[itemize]{leftmargin=1.65em,topsep=0.25em,itemsep=0.14em,parsep=0em}
\setlist[enumerate]{leftmargin=1.8em,topsep=0.25em,itemsep=0.38em,parsep=0em}
\pagestyle{fancy}
\fancyhf{}
\fancyhead[L]{\small Независимые испытания и формула Бернулли}
\fancyhead[R]{\small Екатерина Гнедкова}
\fancyfoot[C]{\small\thepage}
\renewcommand{\headrulewidth}{0.32pt}
\renewcommand{\headrule}{\hbox to\headwidth{\color{accent}\leaders\hrule height \headrulewidth\hfill}}
\renewcommand{\footrulewidth}{0pt}
\newcommand{\edsection}[1]{\vspace{0.45em}\refstepcounter{section}\par\noindent{\color{darkaccent}\fontsize{16}{20}\selectfont \thesection.\ #1}\par\vspace{0.35em}\noindent{\color{accent}\rule{\linewidth}{0.55pt}}\par\vspace{0.35em}}
\newcommand{\minititle}[1]{\par\vspace{0.28em}{\color{darkaccent}\large #1}\par\vspace{0.22em}}
\newcommand{\prob}{\mathsf{P}}
\newcommand{\comb}[2]{\binom{#1}{#2}}
\begin{document}
\thispagestyle{empty}
\begingroup
\centering
\vspace*{0.12\textheight}
{\color{darkaccent}\fontsize{17}{21}\selectfont МАТЕМАТИКА\par}
\vspace{2.15cm}
{\fontsize{25}{31}\selectfont Чек-лист\par}
\vspace{1.15cm}
{\color{accent}\fontsize{23}{29}\selectfont Теория вероятностей:\par}
\vspace{0.1cm}
{\color{accent}\fontsize{21}{27}\selectfont независимые испытания\par}
{\color{accent}\fontsize{21}{27}\selectfont и формула Бернулли\par}
\vspace{1.55cm}
{\fontsize{13}{18}\selectfont Подготовка к ЕГЭ по профильной математике\par}
\vfill
{\fontsize{14}{18}\selectfont Екатерина Гнедкова\par}
\vspace{0.68cm}
{\fontsize{12}{17}\selectfont 6 ноября 2026 г.\par}
\vspace{1.06cm}
{\small Лёвин Артём Александрович\par}
{\small эксклюзивно для Екатерины Гнедковой\par}
\vspace{1.3cm}
\begin{tikzpicture}[x=1cm,y=1cm]
\draw[accent,line width=0.9pt] (-4,0) .. controls (-2.85,0.65) and (-1.35,-0.55) .. (0,0) .. controls (1.40,0.63) and (2.85,-0.54) .. (4,0);
\fill[accent] (0,0) circle (1.45pt);
\end{tikzpicture}
\vspace*{0.45cm}
\endgroup
\clearpage
\setcounter{page}{1}

\edsection{Что нужно уметь}
После повторения темы Вы сможете выбирать способ подсчёта вероятности, различать \emph{ровно}, \emph{хотя бы}, \emph{ни разу}, учитывать заданный порядок и использовать формулу Бернулли.

\edsection{Равновероятные исходы: броски монеты}
Пусть \(\mathrm{O}\) обозначает орла, \(\mathrm{P}\)~--- решку. При двух бросках правильной монеты все четыре упорядоченных исхода равновероятны:
\[
\mathrm{OO},\qquad \mathrm{OP},\qquad \mathrm{PO},\qquad \mathrm{PP}.
\]
Событию «сначала орёл, затем решка» соответствует один исход, поэтому
\[
\prob(\mathrm{OP})=\frac14.
\]
А событию «орёл выпадет ровно один раз» соответствуют два исхода: \(\mathrm{OP}\) и \(\mathrm{PO}\). Следовательно,
\[
\prob(\text{ровно один орёл})=\frac24=\frac12.
\]
Слова «хотя бы один орёл» допускают три исхода: \(\mathrm{OO},\mathrm{OP},\mathrm{PO}\), значит
\[
\prob(\text{хотя бы один орёл})=\frac34.
\]
\minititle{Три броска}
Всего имеется \(2\cdot2\cdot2=8\) равновероятных последовательностей:
\[
\begin{array}{cccc}
\mathrm{OOO}&\mathrm{OOP}&\mathrm{OPO}&\mathrm{OPP}\\[3pt]
\mathrm{POO}&\mathrm{POP}&\mathrm{PPO}&\mathrm{PPP}.
\end{array}
\]
Например, \(\prob(\mathrm{OPO})=\frac18\); орёл \emph{ровно один раз} встречается в трёх исходах, а \emph{более одного раза}~--- в четырёх.
\[
\prob(\text{ровно один орёл})=\frac38,\qquad
\prob(\text{более одного орла})=\frac48=\frac12.
\]
\textcolor{darkaccent}{\textbf{Главный ориентир:}} в серии из \(n\) бросков правильной монеты число равновероятных последовательностей равно \(2^n\).

\clearpage
\edsection{Когда умножать, когда складывать}
\textcolor{darkaccent}{\textbf{«И» (одновременное выполнение условий).}} Для \emph{независимых} событий \(A\) и \(B\):
\[
\prob(A\text{ и }B)=\prob(A)\cdot\prob(B).
\]
\textcolor{darkaccent}{\textbf{«Или» (выбор из вариантов).}} Вероятности складываются напрямую, если варианты \emph{несовместны}: они не могут наступить одновременно. Если события пересекаются, необходимо отдельно учесть их общую часть.

\edsection{Противоположное событие}
Событие \(\overline A\) означает, что \(A\) не произошло. Всегда выполняется
\[
\boxed{\prob(A)=1-\prob(\overline A)}.
\]
Это особенно удобно для формулировки \emph{«хотя бы один»}: её противоположность~--- \emph{«ни одного»}.

\minititle{Пример. Две кофемашины}
Каждая машина независимо выходит из строя с вероятностью \(0{,}2\) и работает с вероятностью \(0{,}8\). Обозначим \(+\)~--- работает, \(-\)~--- сломалась.
\begin{center}
\begin{tabular}{@{}cccl@{}}
\toprule
Первая & Вторая & Вероятность & Событие \\
\midrule
\(+\) & \(+\) & \(0{,}8\cdot0{,}8=0{,}64\) & обе работают \\
\(+\) & \(-\) & \(0{,}8\cdot0{,}2=0{,}16\) & сломалась вторая \\
\(-\) & \(+\) & \(0{,}2\cdot0{,}8=0{,}16\) & сломалась первая \\
\(-\) & \(-\) & \(0{,}2\cdot0{,}2=0{,}04\) & обе сломались \\
\bottomrule
\end{tabular}
\end{center}
Сумма четырёх вероятностей равна единице. Вероятность того, что \emph{хотя бы одна} машина сломается, удобнее считать через противоположное событие:
\[
\prob(\text{хотя бы одна сломается})=1-0{,}8^2=0{,}36.
\]
А хотя бы одна \emph{останется в строю} с вероятностью \(1-0{,}2^2=0{,}96\).

\edsection{Повторные независимые испытания}
Пусть каждое из \(n\) независимых испытаний заканчивается успехом с одной и той же вероятностью \(p\), а неудачей~--- с вероятностью \(q=1-p\).
\[
\begin{aligned}
\prob(\text{все успешны})&=p^n,\\
\prob(\text{ни одного успеха})&=q^n,\\
\prob(\text{хотя бы один успех})&=1-q^n,\\
\prob(\text{хотя бы одна неудача})&=1-p^n.
\end{aligned}
\]

\clearpage
\edsection{Заданные позиции и произвольные позиции}
Биатлонист стреляет по пяти мишеням. Вероятность попадания в каждую равна \(p=0{,}8\), промаха~--- \(q=0{,}2\). Считаем выстрелы независимыми.

\minititle{Первые две мишени поражены, остальные три~--- нет}
Здесь задана \emph{конкретная последовательность} успехов и неудач:
\[
\underbrace{\text{попал, попал}}_{2}\,,\ \underbrace{\text{мимо, мимо, мимо}}_{3}.
\]
Её вероятность:
\[
\prob=p^2q^3=0{,}8^2\cdot0{,}2^3=0{,}00512.
\]
\minititle{Ровно две мишени из пяти поражены}
Теперь попадания могут стоять \emph{на любых двух позициях}. Таких расположений
\[
\comb52=\frac{5!}{2!\,3!}=10.
\]
Вероятность каждого расположения та же, поэтому
\[
\prob=10\cdot0{,}8^2\cdot0{,}2^3=0{,}0512.
\]
Число \(10\) учитывает выбор позиций; при фиксированном порядке этот множитель отсутствует.

\edsection{Формула Бернулли}
Когда проводятся \(n\) \emph{независимых} испытаний с постоянной вероятностью успеха \(p\), вероятность получить \emph{ровно \(k\) успехов} равна
\[
\boxed{\prob_n(k)=\comb nk\,p^k(1-p)^{n-k}},\qquad 0\le k\le n.
\]
Здесь \(n\)~--- число испытаний, \(k\)~--- нужное число успехов,
\(\displaystyle\comb nk=\frac{n!}{k!(n-k)!}\)~--- число способов выбрать позиции успехов; \(0!=1\).

\minititle{Пример. Пять бросков кубика}
Найдём вероятность того, что пятёрка выпадет \emph{ровно три раза}. Для одного броска \(p=\frac16\), \(q=\frac56\); порядок успешных бросков не задан:
\[
\prob_5(3)=\comb53\left(\frac16\right)^3\left(\frac56\right)^2
=10\cdot\frac{25}{6^5}=\frac{125}{3888}.
\]
Если бы было сказано «пятёрка выпала \emph{в первых трёх} бросках и не выпала в остальных двух», ответ был бы \(\frac{25}{6^5}\), без множителя \(10\).

\minititle{Перед записью ответа}
Проверьте, независимы ли испытания, постоянна ли вероятность успеха и указан ли порядок. «Ровно $k$» и «хотя бы $k$» задают разные события.

\clearpage
\thispagestyle{fancy}
\begin{center}
{\color{darkaccent}\fontsize{17}{21}\selectfont Алгоритм выбора метода решения}\par
\vspace{0.5cm}
\end{center}
\begin{center}
\begin{tikzpicture}[
>=Latex, thick,
start/.style={ellipse,draw=accent,fill=white,align=center,text width=3.5cm,minimum height=0.9cm,font=\small},
dec/.style={diamond,draw=accent,aspect=2.1,align=center,text width=3.3cm,inner sep=1pt,font=\small},
act/.style={rectangle,rounded corners=2pt,draw=accent,align=center,text width=4.2cm,minimum height=1.1cm,font=\small},
arr/.style={->,line width=0.85pt,draw=darkaccent}]
\node[start] (s) at (-1.55,0) {Сформулируйте событие};
\node[dec] (d1) at (-1.55,-2.10) {«Хотя бы один»?};
\node[act] (a1) at (4.20,-2.10) {Найдите вероятность «ни одного» и вычтите из единицы};
\node[dec] (d2) at (-1.55,-5.12) {«Ровно $k$» из $n$,\newline позиции не заданы?};
\node[act] (a2) at (4.20,-5.12) {Примените формулу Бернулли};
\node[dec] (d3) at (-1.55,-8.18) {Задан конкретный\newline порядок исходов?};
\node[act] (a3) at (4.20,-8.18) {Перемножьте вероятности отдельных исходов};
\node[dec] (d4) at (-1.55,-11.22) {Мало исходов,\newline все равновероятны?};
\node[act] (a4) at (4.20,-11.22) {Составьте таблицу исходов и подсчитайте долю подходящих};
\node[act] (a5) at (-1.55,-14.02) {Разбейте событие на несовместные случаи и сложите их вероятности};
\draw[arr] (s)--(d1);
\draw[arr] (d1.east)--node[above,font=\scriptsize]{да}(a1.west);
\draw[arr] (d1.south)--node[left,font=\scriptsize]{нет}(d2.north);
\draw[arr] (d2.east)--node[above,font=\scriptsize]{да}(a2.west);
\draw[arr] (d2.south)--node[left,font=\scriptsize]{нет}(d3.north);
\draw[arr] (d3.east)--node[above,font=\scriptsize]{да}(a3.west);
\draw[arr] (d3.south)--node[left,font=\scriptsize]{нет}(d4.north);
\draw[arr] (d4.east)--node[above,font=\scriptsize]{да}(a4.west);
\draw[arr] (d4.south)--node[left,font=\scriptsize]{нет}(a5.north);
\end{tikzpicture}
\end{center}

\clearpage
\begin{center}
{\color{darkaccent}\fontsize{18}{23}\selectfont Тренировочный блок}\par
\vspace{0.15cm}
{\small Путь А \quad\textbullet\quad 10 заданий}\par
\end{center}
\vspace{0.3em}
\textcolor{darkaccent}{\textbf{Указание.}} Если явно не оговорено иное, испытания независимы. Результат можно записать обыкновенной или десятичной дробью. Выберите для каждого задания удобный метод.
\minititle{Базовый уровень}
\begin{enumerate}
\item Правильную монету бросают дважды. Найдите вероятность последовательности \(\mathrm{PO}\) (сначала решка, затем орёл).
\item Правильную монету бросают дважды. Найдите вероятность того, что орёл выпадет ровно один раз.
\item Правильную монету бросают трижды. Найдите вероятность того, что хотя бы один раз выпадет орёл.
\end{enumerate}
\minititle{Уровень выше среднего}
\begin{enumerate}[start=4]
\item Правильную монету бросают четыре раза. Найдите вероятность того, что орёл выпадет более одного раза.
\item Две кофемашины независимо ломаются в течение дня с вероятностью \(0{,}15\) каждая. Какова вероятность того, что хотя бы одна выйдет из строя?
\item Вероятность попадания при каждом из четырёх независимых выстрелов равна \(0{,}7\). Найдите вероятность того, что стрелок попадёт в первые две мишени и промахнётся по последним двум.
\item Игральный кубик бросают шесть раз. Найдите вероятность того, что шестёрка выпадет хотя бы однажды. Ответ запишите в виде выражения, затем вычислите дробь.
\end{enumerate}
\minititle{Сложный уровень}
\begin{enumerate}[start=8]
\item Вероятность попадания при каждом из четырёх независимых выстрелов равна \(0{,}75\). Найдите вероятность ровно двух попаданий.
\item Игральный кубик бросают пять раз. Найдите вероятность того, что шестёрка выпадет ровно дважды.
\item Каждое из шести независимых устройств выходит из строя с вероятностью \(0{,}1\). Найдите вероятность того, что сломаются хотя бы два устройства.
\end{enumerate}

\clearpage
\begin{center}
{\color{darkaccent}\fontsize{18}{23}\selectfont Ответы}\par
\vspace{0.2cm}
{\small Тренировочный блок \(\cdot\) Путь А}\par
\end{center}
\vspace{0.6cm}
\begin{center}
\renewcommand{\arraystretch}{1.55}
\begin{tabular}{@{}p{12mm}p{0.76\linewidth}@{}}
\toprule
№ & Ответ \\
\midrule
1 & \(\dfrac14\) \\
2 & \(\dfrac12\) \\
3 & \(\dfrac78\) \\
4 & \(\dfrac{11}{16}\) \\
5 & \(1-0{,}85^2=0{,}2775\) \\
6 & \(0{,}7^2\cdot0{,}3^2=0{,}0441\) \\
7 & \(1-\left(\dfrac56\right)^6=\dfrac{31031}{46656}\) \\
8 & \(\comb42\cdot0{,}75^2\cdot0{,}25^2=\dfrac{27}{128}\) \\
9 & \(\comb52\left(\dfrac16\right)^2\left(\dfrac56\right)^3=\dfrac{625}{3888}\) \\
10 & \(1-0{,}9^6-6\cdot0{,}1\cdot0{,}9^5=0{,}114265\) \\
\bottomrule
\end{tabular}
\end{center}
\vspace{1.1cm}
\minititle{Краткая памятка}
\begin{center}
\begin{tabularx}{\linewidth}{@{}p{0.40\linewidth}X@{}}
\toprule
Сигнал в условии & Рекомендуемый приём \\
\midrule
«Хотя бы один» & Противоположное событие \\
«Все» / «ни одного» & Произведение одинаковых вероятностей \\
«Ровно $k$ из $n$» & Формула Бернулли \\
«Первые два», «последний» & Фиксированная последовательность \\
Два-три броска монеты & Таблица равновероятных исходов \\
\bottomrule
\end{tabularx}
\end{center}
\vfill
\begin{center}
{\small\color{textgray}Лёвин Артём Александрович\quad\textbullet\quad эксклюзивно для Екатерины Гнедковой}
\end{center}
\end{document}
